\Subhead{5.1.1}{Вибратор Герца. Ближняя и дальняя зоны}

Вибратор Герца — идеализированный короткий элемент длины $\ell\ll\lambda$ с одинаковым по фазе и величине током $I_0$ вдоль элемента. Это модель элементарного излучателя; у реального короткого провода ток может быть треугольным и эффективная длина отличается. Дипольный момент связан с током $I_0\ell=-\ii\omega p_0$ при $\e^{-\ii\omega t}$.

Поле диполя содержит члены $1/r^3$, $1/r^2$, $1/r$: квазистатический, индукционный и излучательный. Например,
\begin{equation}
 \bm E=\frac{\e^{\ii kr}}{4\pi\varepsilon_0}
 \left\{\frac{k^2}{r}(\bm n\times\bm p_0)\times\bm n+
 \left(\frac1{r^3}-\frac{\ii k}{r^2}\right)[3\bm n(\bm n\cdot\bm p_0)-\bm p_0]\right\}.
\end{equation}
В дальней зоне $kr\gg1$ остаются поперечные $E_\theta,H_\varphi$, связанные $E_\theta=Z_{\rm vac}H_\varphi$:
\begin{equation}
 E_\theta=-\ii\frac{Z_{\rm vac}kI_0\ell}{4\pi r}\sin\theta\,\e^{\ii kr},
 \qquad \frac{\dd P}{\dd\Omega}=\frac{Z_{\rm vac}k^2|I_0|^2\ell^2}{32\pi^2}\sin^2\theta.
\end{equation}
Интегрирование $\sin^2\theta$ по телесному углу даёт $P=Z_{\rm vac}k^2|I_0|^2\ell^2/(12\pi)$ и сопротивление излучения
\begin{equation}
 R_{\rm rad}=\frac{2P}{|I_0|^2}=80\pi^2(\ell/\lambda)^2.
\end{equation}
Амплитуда тока здесь пиковая. Ближняя зона содержит значительный реактивный обмен, и отношение $E/H$ не равно плоско-волновому импедансу. Для антенны большого размера $D$ дальняя зона требует также малой квадратичной фазовой ошибки: обычно $r\gg2D^2/\lambda$, а не только $r\gg\lambda$. Диаграмма диполя имеет нули вдоль оси и максимум поперёк.
